\documentclass{article}

\usepackage{amsmath,amssymb}
\usepackage{xurl}
\usepackage[hidelinks]{hyperref}
\setlength{\emergencystretch}{2em}

\input{command}

\title{Scope of Ha's $A_\gamma$ Witness for Choi-Type Maps}
\author{}
\date{2026-08-24}

\begin{document}
\maketitle
\gapnote{scope-restriction}{wip}

\section{Source theorem and witness}

Let $d\geq3$, let $1\leq n\leq d-2$, and let $D(X)$ denote the diagonal
part of $X\in\MD$. For the Weyl operator $U_{k0}$, Wolf's Example~3.1 defines
the Choi-type map
\begin{align}
    T_C(X)
        &=(d-n)D(X)-X
          +\sum_{k=1}^{n}D(U_{k0}XU_{k0}^{\dagger}).
    \label{eq:ha98_choi_type_witness_scope_choi_type_map}
\end{align}
The source identifies this formula as Equation~(3.20) and states that $T_C$ is
positive and indecomposable \cite[Example~3.1]{Wolf2012Quantum}. The local
source is
\path{Notes/WolfNoteTexSource/ch03_positive_not_completely.tex}, lines
357--365. Wolf gives neither a proof nor a citation in that paragraph.

Ha proves the stronger atomicity assertion for the same family
\cite[Theorem~2.1, pp.~593--595]{Ha1998AtomicPositiveMaps}. Fix
$0<\gamma<1$. Ha takes all $3^d$-th roots of unity, constructs 2-simple
vectors $z_{r,i}$, and defines
\begin{align}
    A_r
        &=\frac{1}{3^d}\sum_{i=1}^{3^d}\ketbra{z_{r,i}}{z_{r,i}},
    \label{eq:ha98_choi_type_witness_scope_root_average_ar}\\
    A_\gamma
        &=\frac{1}{d}\sum_{r=1}^{d}A_r.
    \label{eq:ha98_choi_type_witness_scope_root_average_agamma}
\end{align}
Let $V_2$ be the cone generated by projectors onto 2-simple vectors. The
construction in
(\ref{eq:ha98_choi_type_witness_scope_root_average_ar})--%
(\ref{eq:ha98_choi_type_witness_scope_root_average_agamma}) gives
$A_\gamma\in V_2$. Ha derives the entries of $A_\gamma$ from the root
power sum and the exponent-collision identity in Equations~(2.1)--(2.3) of
Ref.~\cite{Ha1998AtomicPositiveMaps}.

Write $A^T$ for Ha's block transpose, namely transposition in the second tensor
factor. With cyclic suffixes, Ha defines
\begin{align}
    u_i
        &=\frac{\gamma}{\sqrt d}e_{i+1}\otimes e_i
          +\frac{1}{\sqrt d\,\gamma}e_i\otimes e_{i+1},
    \label{eq:ha98_choi_type_witness_scope_ui}\\
    v_i
        &=\sqrt{\frac{d-1}{d}}
          (e_{i+1}\otimes e_i+e_i\otimes e_{i+1}),
    \label{eq:ha98_choi_type_witness_scope_vi}\\
    \alpha_i
        &=e_i\otimes e_i.
    \label{eq:ha98_choi_type_witness_scope_alpha_i}
\end{align}
The source also defines vectors $\beta_{1j}$ and $\beta_{ij}$ for
non-neighbouring pairs. Its calculation on p.~595 is
\begin{align}
    A_\gamma^T
        &=\sum_{i=1}^{d}
          \bigl(          \ketbra{u_i}{u_i}
              +\ketbra{v_i}{v_i}
              +\ketbra{\alpha_i}{\alpha_i}
          \bigr)
          +\sum_{j=3}^{d-1}\ketbra{\beta_{1j}}{\beta_{1j}}
          +\sum_{i=2}^{d-2}\sum_{j=i+2}^{d}
            \ketbra{\beta_{ij}}{\beta_{ij}}.
    \label{eq:ha98_choi_type_witness_scope_block_transpose_decomposition}
\end{align}
Every summand in
(\ref{eq:ha98_choi_type_witness_scope_block_transpose_decomposition}) is
generated by a 2-simple vector, so $A_\gamma^T\in V_2$. Ha finally computes
\begin{align}
    \langle A_\gamma,T_C\rangle
        &=\gamma^2-1
        <0.
    \label{eq:ha98_choi_type_witness_scope_negative_pairing}
\end{align}

\section{Formalized construction}

Ha uses the one-based integers
$m_k=\frac{3}{2}(3^{k-1}-1)$ and the base vector $a_{i,1}$. The formalization
uses zero-based indices, so the corresponding definitions are
$m_k=\frac{3}{2}(3^k-1)$ and $a_{i,0}$. The same convention defines the root
family, the vectors $a_{i,k}$, $c_r$, $b_{r,i,j}$, and $z_{r,i}$, and the
averages in
(\ref{eq:ha98_choi_type_witness_scope_root_average_ar})--%
(\ref{eq:ha98_choi_type_witness_scope_root_average_agamma}). Cyclic
coordinates lie in $\mathbb Z/d\mathbb Z$, consistently with Wolf's
Equation~(2.24) \cite{Wolf2012Quantum} and Ha's rotation
$Se_i=e_{i+1}$ \cite{Ha1998AtomicPositiveMaps}.

Fix $r$ and $i$. Every column of $z_{r,i}$ except column $r$ is a scalar
multiple of $a_{i,0}$, while column $r$ is $c_r\circ a_{i,r}$. Its coefficient
matrix therefore has rank at most two. Averaging the associated rank-one
projectors proves
\begin{align}
    A_\gamma
        &\in V_2,
    \label{eq:ha98_choi_type_witness_scope_agamma_vtwo}\\
    A_\gamma
        &\geq0.
    \label{eq:ha98_choi_type_witness_scope_agamma_positive}
\end{align}

The formalization also proves Ha's root-average entry formula, Equation~(2.3)
of Ref.~\cite{Ha1998AtomicPositiveMaps}. It defines the vectors in
(\ref{eq:ha98_choi_type_witness_scope_ui})--%
(\ref{eq:ha98_choi_type_witness_scope_alpha_i}) and the vectors
$\beta_{i,j}$, with Ha's two beta ranges translated to zero-based indices.
Let $\mathcal B_d$ denote the union of those ranges, and let $T_2$ denote right
partial transpose. The resulting identity and cone membership are
\begin{align}
    A_\gamma^{T_2}
        &=\sum_{i=0}^{d-1}
          \bigl(          \ketbra{u_i}{u_i}
              +\ketbra{v_i}{v_i}
              +\ketbra{\alpha_i}{\alpha_i}
          \bigr)
          +\sum_{(i,j)\in\mathcal B_d}
            \ketbra{\beta_{i,j}}{\beta_{i,j}},
    \label{eq:ha98_choi_type_witness_scope_formal_block_transpose}\\
    A_\gamma^{T_2}
        &\in V_2.
    \label{eq:ha98_choi_type_witness_scope_formal_block_transpose_vtwo}
\end{align}
For $d=3$, the beta range is empty, as it is in the source. Thus the
formalization proves both source-ordered $V_2$ decompositions on pp.~594--595
of Ref.~\cite{Ha1998AtomicPositiveMaps}.\footnote{Formal declarations:
    \leanid{Matrix.haTwoSimpleVector\_hasSchmidtRankLE\_two},
    \leanid{Matrix.haArGamma\_apply},
    \leanid{Matrix.haAGamma\_hasSchmidtNumberLE\_two}, and
    \leanid{Matrix.haAGamma\_posSemidef} in
    \doclink{QICLean/Channel/ChoiTypeMap/HaTwoSimpleWitness};
    \leanid{Matrix.partialTransposeRight\_haAGamma\_eq\_haBlockTransposeDecomposition}
    and
    \leanid{Matrix.partialTransposeRight\_haAGamma\_hasSchmidtNumberLE\_two}
    in \doclink{QICLean/Channel/ChoiTypeMap/HaBlockTranspose}.}

\section{Factor swap, pairing, and remaining scope}

Let $\sigma$ swap the tensor factors, and define
\begin{align}
    J_d
        &=\sum_{i,j}e_{ij}\otimes e_{ij}.
    \label{eq:ha98_choi_type_witness_scope_jd}
\end{align}
The Eom--Kye duality convention gives
\begin{align}
    \langle A,\Phi\rangle
        &=\bigl\langle
            (\Phi\otimes\operatorname{id})(A^\sigma),J_d
          \bigr\rangle.
    \label{eq:ha98_choi_type_witness_scope_eom_kye_orientation}
\end{align}
See equations~(12)--(13), Theorem~3.3, and the concluding shuffled pairing in
Ref.~\cite[pp.~137--139]{EomKye2000Duality}. Consequently, the state tested by
the ampliation and its left partial transpose are
\begin{align}
    \rho
        &=A_\gamma^\sigma,
    \label{eq:ha98_choi_type_witness_scope_witness_state}\\
    \rho^{T_1}
        &=(A_\gamma^{T_2})^\sigma.
    \label{eq:ha98_choi_type_witness_scope_partial_transpose_transport}
\end{align}
The decompositions in
(\ref{eq:ha98_choi_type_witness_scope_agamma_vtwo}) and
(\ref{eq:ha98_choi_type_witness_scope_formal_block_transpose_vtwo}) transport
through the factor swap. The formalization therefore proves
$\rho\in V_2$ and $\rho^{T_1}\in V_2$.

The formal development uses the normalized maximally entangled vector
\begin{align}
    \ket{\Omega}
        &=\frac{1}{\sqrt d}\sum_i e_i\otimes e_i,
    \label{eq:ha98_choi_type_witness_scope_normalized_omega}\\
    J_d
        &=d\ketbra{\Omega}{\Omega}.
    \label{eq:ha98_choi_type_witness_scope_omega_normalization}
\end{align}
It retains the positive factor $d$. For the full range $d\geq3$,
$1\leq n\leq d-2$, and $\gamma>0$, the formalized identity is
\begin{align}
    \langle A_\gamma,T_C\rangle
        &=d\,\bra{\Omega}
          (T_C\otimes\operatorname{id})(A_\gamma^\sigma)
          \ket{\Omega},
    \label{eq:ha98_choi_type_witness_scope_pairing_quadratic_form}\\
    \langle A_\gamma,T_C\rangle
        &=\gamma^2-1.
    \label{eq:ha98_choi_type_witness_scope_formal_negative_pairing}
\end{align}
Thus $0<\gamma<1$ gives the strict negativity in
(\ref{eq:ha98_choi_type_witness_scope_negative_pairing}).\footnote{The factor
swap, partial-transpose transport, pairing identity, cyclic-coordinate
calculation, and normalized quadratic form are in
\doclink{QICLean/Channel/ChoiTypeMap/HaNegativePairing}; the principal
declarations are
\leanid{Matrix.partialTransposeLeft\_tensorFactorSwap\_haAGamma\_hasSchmidtNumberLE\_two},
\leanid{Matrix.eomKyePairing\_eq\_omegaVec\_quadraticForm\_factorSwap},
\leanid{Matrix.eomKyePairing\_haAGamma\_choiTypeMap}, and
\leanid{Matrix.choiTypeMapFin\_haAGamma\_omegaVec\_quadraticForm\_re\_neg}.}

This proves the source algebra of the witness. Positivity throughout the source
range is also formalized, and the former middle-range gap is resolved in
\ghissue{19}.\footnote{Formal declaration:
\leanid{Matrix.choiTypeMap\_isPositiveMap}.}
A theorem combining positivity with the negative witness to conclude
indecomposability remains to be stated; \ghissue{18} tracks this final step.
The present formalization does not assert atomicity.

\section{Classification}

\emph{Scope restriction.} The formalization covers the exact root-average
construction, its entry formula, both $V_2$ decompositions on pp.~594--595 of
Ref.~\cite{Ha1998AtomicPositiveMaps}, their tensor-factor transport, the
Eom--Kye orientation, the normalization factor $d$, and the value
$\gamma^2-1$ for every $d\geq3$ and $1\leq n\leq d-2$. It also proves
positivity in this range. The final indecomposability theorem remains open.
No atomicity conclusion or alternative Choi-matrix or partial-transpose ansatz
is asserted.

\bibliographystyle{alpha}
\bibliography{references}

\end{document}
